\documentclass[a4paper]{article}
\usepackage[english]{babel}
\usepackage[utf8]{inputenc}
\usepackage[colorinlistoftodos]{todonotes}
\usepackage{amsmath,amsthm,geometry,graphics,amssymb,csquotes,hyperref,float,amsfonts,natbib}
\geometry{verbose,tmargin=3cm,bmargin=3cm,lmargin=3cm,rmargin=3cm}
\theoremstyle{plain}
\newtheorem{prop}{Proposition}
\newtheorem{theorem}{Theorem}
\newtheorem{corollary}{Corollary}[section]
\newtheorem{lemma}{Lemma}[section]
\newtheorem{definition}{Definition}[section]
\newtheorem*{remark}{Remark}
\renewcommand\qedsymbol{$\blacksquare$}
%%%%%%%%%%%%%%%%%%%%%%%%%%%%%%%%%%%%%%%%%%%%%%%%%%%%%%%%%%%%%%%%%%%%%%%%%%%%%%%%%%%%%%%%%%%%
%%%%%%%%%%%%%%%%%%%%%%%%%%%%%%%%%%%%%%%%%%%%%%%%%%%%%%%%%%%%%%%%%%%%%%%%%%%%%%%%%%%%%%%%%%%%
%%%%%%%%%%%%%%%%%%%%%%%%%%%%%%%%%%%%%%%%%%%%%%%%%%%%%%%%%%%%%%%%%%%%%%%%%%%%%%%%%%%%%%%%%%%%
\title{Note on Asymptotic Theory}
\author{Haixi Li}
\date{\today}
%%%%%%%%%%%%%%%%%%%%%%%%%%%%%%%%%%%%%%%%%%%%%%%%%%%%%%%%%%%%%%%%%%%%%%%%%%%%%%%%%%%%%%%%%%%
%%%%%%%%%%%%%%%%%%%%%%%%%%%%%%%%%%%%%%%%%%%%%%%%%%%%%%%%%%%%%%%%%%%%%%%%%%%%%%%%%%%%%%%%%%%
%%%%%%%%%%%%%%%%%%%%%%%%%%%%%%%%%%%%%%%%%%%%%%%%%%%%%%%%%%%%%%%%%%%%%%%%%%%%%%%%%%%%%%%%%%%
\begin{document}
\maketitle
Notation in the paper follows \cite{Tartakovsky2005}.
\section{Problem formulation}
The observed sequence of random variables $X_{1},X_{2},...$ are i.i.d., untile a change occurs, the observations are again i.i.d., but with another distribution.
General guideline for math notations:
\begin{itemize}
    \item use lower case Greek letter for parameters, $\beta ,\alpha $
    \item use capital letters for random variables, $X$,$Y$, sometime for matrix
    \item use lower case letter for realization of random variables, $x$,$y$, or individual observation.
    \item use $\mathbb{R}$ denote space and $\mathcal{F}$ to denote $\sigma$-algebra on the space.
\end{itemize}
\section{Measure}
A measure is a set function, a mapping which associates a real number with a set.
\cite{Davidson1994}
\section{Expectation}
\begin{corollary}
Let $g:\mathbb{R}\mapsto \mathbb{R}$ be a function with that property that $x\geq a$ implies $g(x)\geq g(a)>0$, for a given constant $a$. Then 
\[
P(X\geq a) \leq \frac{E(g(X))}{g(a)}.
\]
\end{corollary}
\begin{proof}
\begin{eqnarray*}
g(a)P(X\geq a) &=& g(a)\int_{x\geq a} dF(x)\\
               &\leq& \int_{x\geq a}g(x)dF(x) \\
               &\leq& E(g(X).
\end{eqnarray*}
\end{proof}
\section*{Appendix}
\begin{table}[H]
	\footnotesize
	\centering
	\begin{tabular}{c l}
		\hline
		\hline
		Notations             & Description \\
		\hline
		$\lambda$	& structural breaking time \\
		$\tau$	& stopping time\\
		$\Omega$& sample space\\
		$\omega \in \Omega$& a point in the sample space\\
		$\mathbb{R}$ & real number space\\
		$\mathbb{Q}$ & rational number space\\
		$(\Omega,\mathcal{F},\mu)$ & measurable space\\
		\hline
		\hline
	\end{tabular}
	\protect\caption{\label{notations}Table of Notations}
	\end{table}
\bibliographystyle{elsarticle-harv}
\bibliography{/home/ethorondor/Documents/Mendeley/library.bib}
\end{document}